\section{Theory (fixed-node / root guarantees)}
\label{sec:theory}

This section states the inferential guarantees we use in the paper, in the
strongest defensible form. Proofs are provided in the appendices.

\subsection{Setup and assumptions}
\label{sec:theory-setup}

We state results for a \emph{fixed} node $t$ with index set $I_t$ and data
$(X_t, Y_t)$. Let $F_t$ be the tested feature set at node $t$, with
$m_t := |F_t|$.

Our guarantees are finite-sample and fixed-node: they are cleanest at the root.
We keep all assumptions explicit and aligned across results; see the assumptions
ledger in \cref{app:assumptions} and the compact checklist in
\cref{app:assumption-checklist}. The core requirements are conditional
exchangeability under the null, a label-independent tested family and resampling
budget, fixed-$B$ p-values (no optional stopping), and a conservative or
randomized tie convention.

\subsection{Validity of the +1 Monte Carlo permutation p-value}
\label{sec:theory-plusone}

\begin{theorem}[+1 permutation p-values are super-uniform]
\label{them:plusone-superuniform}
Assume A0.1--A0.5 (\cref{app:assumptions}). In particular, under the null
hypothesis being tested, the permutation statistics $(T_0, T_1, \dots, T_B)$ are
exchangeable (conditional on $(X_t, U)$). Define the +1 Monte Carlo p-value
%
\begin{equation}
  p := \frac{1 + \sum_{b=1}^B \ind\{T_b \ge T_0\}}{B+1}.
\end{equation}
%
Then for all $\alpha \in [0,1]$, $\PP(p \le \alpha) \le \alpha$.
\end{theorem}

\subsection{Stage A control via Bonferroni}
\label{sec:theory-stageA}

\begin{proposition}[Stage A global-null bound at a fixed node]
\label{prop:stageA-global-null}
Assume A0.1--A0.5 (\cref{app:assumptions}).
Assume $H^{\mathrm{sel}}_{t,j}$ is true for all $j \in F_t$ and each $p_{t,j}$ is
super-uniform. If Stage~A uses Bonferroni and rejects when
$\min_{j \in F_t} p_{t,j} \le \alpha_{\mathrm{sel}}/m_t$, then
%
\begin{equation}
  \PP\!\left(\exists j \in F_t : p_{t,j} \le \alpha_{\mathrm{sel}}/m_t\right) \le \alpha_{\mathrm{sel}}.
\end{equation}
\end{proposition}

\begin{proposition}[Per-feature false selection bound (cardinality-free)]
\label{prop:per-feature-bound}
Assume A0.1--A0.5 (\cref{app:assumptions}).
Fix a feature $j \in F_t$ whose Stage~A null $H^{\mathrm{sel}}_{t,j}$ is true, so
$p_{t,j}$ is super-uniform. If Stage~A uses Bonferroni at level
$\alpha_{\mathrm{sel}}/m_t$, then
%
\begin{equation}
  \PP(\text{node $t$ splits on feature $j$}) \le \alpha_{\mathrm{sel}}/m_t.
\end{equation}
%
In particular, the bound does not depend on the number of candidate thresholds
available for feature $j$.
\end{proposition}

\subsection{A root-level global statement}
\label{sec:theory-root}

\begin{proposition}[Any split implies root Stage A rejection]
\label{prop:any-split-implies-root-rejection}
Assume A0.1--A0.5 at the root (\cref{app:assumptions}).
Consider a tree grown with Stage~A screening at the root in fixed-$B$ mode. The
fitted tree can only contain any internal split if the root passes Stage~A.
Consequently, under a complete global null over tested root features,
%
\begin{equation}
  \PP(\text{the fitted tree has at least one internal split}) \le \alpha_{\mathrm{sel}}.
\end{equation}
\end{proposition}

\subsection{Multi-selector mode}
\label{sec:theory-multiselector}

\begin{proposition}[Multi-selector max-statistic validity]
\label{prop:multiselector-validity}
Assume A0.1--A0.5 (\cref{app:assumptions}). Fix a feature at a fixed node and
let $T^{(1)}_b, \dots, T^{(K)}_b$ denote $K$ selector statistics computed on
permutation $b$, and define $M_b := \max_k |T^{(k)}_b|$. If the joint vector of
all selector statistics is exchangeable across permutations under the null,
then the +1 p-value computed from $(M_b)$ is super-uniform.
\end{proposition}

\noindent See \cref{app:multiselector} for details and a discussion of
scale/normalization considerations when mixing selectors.
